% Wspólna podstawa geometryczna dla rozdmuchania punktu.
% Wariant opisu przez normalną i przez konormalną: obydwa zachowane.
% Historyczne etykiety występują tylko przy pierwszym włączeniu.
\ifcsname biblioteka@rozdmuchanie@normalna@loaded\endcsname
Wiązkę normalną dzielnika wyjątkowego obliczono wcześniej
w lemacie~\ref{lem:xvii-exceptional}.
\else
\expandafter\gdef\csname biblioteka@rozdmuchanie@normalna@loaded\endcsname{}
\begin{lemat}[Wiązka normalna dzielnika wyjątkowego]
\label{st:ga6-exceptional-normal}\label{lem:xvii-exceptional}
Niech $S$ będzie gładką zespoloną powierzchnią, $p\in S$, a
$\pi:\widetilde S=\operatorname{Bl}_pS\to S$ jej rozdmuchaniem
w punkcie $p$. Wtedy $E=\pi^{-1}(p)\simeq\mathbb P^1_{\mathbb C}$,
$N_{E/\widetilde S}\simeq\mathcal O_{\mathbb P^1}(-1)$ i
$E\cdot E=\deg N_{E/\widetilde S}=-1$.
Ostatnią liczbę rozumiemy jako lokalne samoprzecięcie krzywej,
nawet jeśli $S$ nie jest zwarta.
\end{lemat}
\begin{proof}
Wybierzmy współrzędne $(x,y)$ powierzchni $S$ z początkiem w $p$.
Rozdmuchanie jest lokalnie pododmianą
$\{((x,y),[s:t])\colon xt=ys\}\subset\mathbb C^2\times\mathbb P^1$.
W mapie $s\ne0$ piszemy $v=t/s$, $y=xv$, a w mapie $t\ne0$
piszemy $u=s/t$, $x=uy$. Na części wspólnej zachodzi
$u=v^{-1}$ i $y=vx$. W pierwszej mapie ideał $E$ jest
generowany przez $x$, w drugiej przez $y$.

\emph{Opis przez wiązkę normalną.}
W punkcie $[1:v]\in E$ kierunek normalny do $E$ jest
reprezentowany przez wektor $(x,xv)=x(1,v)$.
Po zmianie mapy wektory te wyznaczają dokładnie
podwiązki prostych rozpiętych przez $(s,t)$ w
$\mathbb P^1\times\mathbb C^2$; jest to tautologiczna
wiązka $\mathcal O_{\mathbb P^1}(-1)$.

\emph{Równoważny opis przez wiązkę konormalną.}
Lokalne generatory $x$ i $y$ wiązki
$\mathcal I_E/\mathcal I_E^2$ spełniają $y=vx$;
odpowiada to wiązce $\mathcal O_{\mathbb P^1}(1)$.
Jej dualna jest więc $N_{E/\widetilde S}\cong\mathcal O(-1)$.
Ponieważ stopień $\mathcal O_{\mathbb P^1}(-1)$ wynosi $-1$,
formuła samoprzecięcia
$E^2=\langle c_1(N_{E/\widetilde S}),[E]\rangle$
daje $E^2=-1$.
\end{proof}
\fi
